\documentclass[11pt]{article}
\usepackage[no-math]{fontspec}
\usepackage[french,english]{babel}
\usepackage{amsmath}
\usepackage{amssymb}
\usepackage{amsthm}
\usepackage{mathrsfs}
\usepackage{tikz-cd}
\usepackage{enumitem}

\frenchsetup{AutoSpacePunctuation=false}

\newcommand{\printedpage}[1]{\marginpar{\footnotesize p.~#1}}
\newcommand{\msp}[1]{\marginpar{\footnotesize [#1]}}
\newcommand{\und}[1]{\underline{\mathrm{#1}}}
\newcommand{\Cat}{\mathrm{Cat}}
\newcommand{\circled}[1]{\tikz[baseline=(c.base)]\node[draw,circle,inner sep=1pt](c){#1};}

\begin{document}

\printedpage{14}
% head: A. Grothendieck : Dérivateurs, Ch. VI, Édition des Universités de Montpellier II et Paris VII
% foot: Version 9 novembre 2001

\msp{page 22}
Here again, the construction is functorial, we get the functor
\[ C : \Cat/S \longrightarrow \und{Parf}\,S \]
which must be composed with the functor in the opposite direction
\[ r : \Cat/S \longleftarrow \und{Parf}\,S. \]
First of all, let us check that $C$, just like $r$, is compatible with the localizers. Let $f : X \longrightarrow Y$ be in $\Cat/S$, suppose $f$ is in $W$, let us prove $C(f)$ is in $W_S$. We have commutativity in
\[
\begin{tikzcd}[row sep=large, column sep=large]
X \arrow[r, "f"] \arrow[d, "j_X"'] & Y \arrow[d, "j_Y"] \\
CX \arrow[r, "Cf"] & CY
\end{tikzcd}
\]
and $j_X$, $j_Y$ are in $W$, so $f$ is in $W$ if and only if $Cf$ is. But in the case $W = W_\infty$ we know that this is equivalent to $Cf$ being in $W_\infty$ fiber by fiber. Admitting \circled{1} p. 15 for $W$, it remains to make sure of a converse in \underline{Parf}, i.e.

\circled{2} Let
\begin{tikzcd}[baseline=(\tikzcdmatrixname-2-2.base), row sep=small, column sep=small]
X \arrow[rr, "f"] \arrow[dr] & & Y \arrow[dl] \\
& S &
\end{tikzcd}
be in $\und{Parf}\,S$ (or merely in $\und{Fibhot}\,S$). If $f$ is in $W$, then so are the $f_s : X_s \longrightarrow Y_s$.

\msp{page 23}
Let us try to prove it, by copying the proof in the case $C = C_\infty$ (XI, th.~2). We can suppose $S$ connected (\footnote[14]{serves no purpose}), using La (i.e. $W \subseteq W_0$) and admitting the

\textbf{Lemma.} \textit{Let $f : X \longrightarrow Y$ be in} Cat. \textit{Then $f \in W$ if and only if $\pi_0(f)$ is bijective, and for every connected component $X_i$ of $X$, if $Y_i$ is the corresponding connected component of $Y$, $f_i : X_i \longrightarrow Y_i$ induced by $f$ is in $W$.}

Let us then use the existence of
\[ S' \longrightarrow S \]
a Hot-fibration with $S'$ Hot-aspheric (\footnote[15]{\textbf{NB} It suffices to have a Hot-fibration $S' \longrightarrow S_0 \hookrightarrow S$, where $S_0$ is the connected component of $S$ containing the element $s \in \mathrm{Ob}\,S$ under consideration.}). We can even suppose that $S' \longrightarrow S$ is perfect. Suppose we know that base change by such perfect fibrations transforms $W$-equivalences into likewise. (That must be true for the $W_i$, but we have seen that it was false for $k$-homological equivalence, cf. I p.~87 ff.) Then $X' \xrightarrow{f'} Y'$ over $S'$ is
\printedpage{15}
% head: A. Grothendieck : Dérivateurs, Ch. VI, Édition des Universités de Montpellier II et Paris VII
% foot: Version 9 novembre 2001
a $W$-equivalence. Now for $s' \in S'$, since $\{s'\} \longrightarrow S'$ is a Hot-equivalence ($S'$ Hot-aspheric), and $X'$, $Y'$ being Hot-fibered over $S$, it follows that $X'_{s'} \longrightarrow X'$, $Y'_{s'} \longrightarrow Y'$ are Hot-equivalences, hence also $X'_{s'} \xrightarrow{f'_{s'}} Y'_{s'}$, as we see from the commutative square
\[
\begin{tikzcd}[row sep=large, column sep=large]
X'_{s'} \arrow[r, "f'_{s'}"] \arrow[d, "\wr"'] & Y'_{s'} \arrow[d, "\wr"'] \\
X' \arrow[r, "{f'}"] & Y'.
\end{tikzcd}
\]

\msp{page 24}
Since $S' \longrightarrow S$ is surjective, this shows that the $f_s : X_s \longrightarrow Y_s$ are in $W$, qed.

Thus, (2) will be a consequence of \circled{2.} Let $p : Y' \longrightarrow Y$ be a Hot-fibration (or even a perfect functor, if that can help \dots). Let $f : X \longrightarrow Y$ be an $S$-arrow (in $\Cat/S$), then $f' : X' \longrightarrow Y'$ by base change. If $f \in W$, $f' \in W'$
\[
\begin{tikzcd}[row sep=large]
X' \arrow[r, "q"] \arrow[d, "{f'}"'] \arrow[dr, phantom, "\scriptstyle\text{cart.}"] & X \arrow[d, "f"] \\
Y' \arrow[r, "p"] & Y
\end{tikzcd}
\qquad\qquad
f \in W \text{ and } p \text{ perfect} \Longrightarrow f' \in W
\]

In any case, we absolutely need \circled{2} to be able to conclude that $C$ is compatible with the localizers.

Let us consider $rC : \Cat/S \longrightarrow \Cat/S$, we have an arrow
\[
\left\{
\begin{array}{l}
\mathrm{id}_{\Cat/S} \xrightarrow{j} rC \\[1.5ex]
j_X : X \longrightarrow CX \quad \text{the canonical arrow } i_{X/S}\ \text{\texttt{[rather} } j_{X/S}\text{\texttt{]}} \text{ (p.~21),}
\end{array}
\right.
\]
and $j_X$ is always in $W$.

Let us consider $Cr : \und{Parf}\,S \longrightarrow \und{Parf}\,S$, i.e. let us apply the construction $C$ to an $X$ perfect over $S$. Then we again have
\[
\left\{
\begin{array}{l}
\mathrm{id}_{\und{Parf}\,S} \xrightarrow{j'} Cr \\[1ex]
j'_X = j_{r(X)} : X = r(X) \xrightarrow{j_X} CX.
\end{array}
\right.
\]
This morphism is in $W$, hence also a $W$-equivalence fiber by fiber, since source and target are perfect, hence Hot-fibered over $S$.

\msp{page 25}
This already proves that the functors $r$ and $C$ induce equivalences quasi-inverse to one another on the localized categories. By the lemma p.~11, this shows that for every full subcategory $\mathcal{F}$ of $\Cat/S$ containing $\und{Parf}\,S$, the inclusion of $\mathcal{F}$ in $\Cat/S$ induces an
\printedpage{16}
% head: A. Grothendieck : Dérivateurs, Ch. VI, Édition des Universités de Montpellier II et Paris VII
% foot: Version 9 novembre 2001
equivalence on the localized categories, taking as localizer in $\mathcal{F}$ the ``coarse'' localizer induced by $\Cat/S$. This applies in particular to the full subcategories
\[ \und{Liss}\,S, \quad \und{Fibhot}\,S, \quad \und{Prop}\,S, \]
but taking in $\und{Liss}\,S$ and $\und{Prop}\,S$ not the fine ``fiber by fiber'' localizer, as I had wrongly announced, but the coarse localizer induced by $\Cat/S$. (For $\und{Fibhot}\,S$, there is no risk of ambiguity). I ought to convince myself that if I take the fine $W$-equivalence (in the case $W = W_\infty$, say), I do \textit{not} find the ``good'' localized categories, i.e. the functor of localized categories, to the localized category of $\Cat/S$, is not an equivalence. This in fact supposes that in $\und{Fib}\,S$ and $\und{Cofib}\,S$ as well, it is the \textit{coarse} equivalence
\msp{page 26}
that should be taken. (In $\und{Fib}_0S$ and $\und{Cofib}_0S$, or $\und{Fibsc}_0S$ and $\und{Cofibsc}_0S$, the two equivalences are identical, since these four categories are contained in $\und{Fibhot}\,S$.) The sequel will show us whether this is indeed so. On the other hand, in $\und{Fibsc}\,S$ and $\und{Cofibsc}\,S$, there is no doubt, according to the general principles of constructing a derivator from an $(\mathcal{M}, W)$, that we must take the fine equivalence, fiber by fiber. In this connection, we must be aware that the vertical functors of the diagram p.~8
\[ \und{Fibsc}\,S \longrightarrow \und{Fib}\,S, \quad \und{Cofibsc}\,S \longrightarrow \und{Cofib}\,S, \]
though essentially surjective and faithful, are not fully faithful. They would be, and even equivalences of categories, if in $\und{Fib}\,S$ resp. $\und{Cofib}\,S$ I confined myself to the $S$-functors that are \textit{cartesian} resp. \textit{cocartesian}. So we can interpret, up to equivalence, $\und{Fibsc}\,S$ and $\und{Cofibsc}\,S$ as $\und{Fib}\,S$ and $\und{Cofib}\,S$, where we would have restricted the morphisms to just those $S$-functors that are \textit{cartesian} resp. \textit{cocartesian.} It would thus seem that by restricting the arrows in this way, but by restricting the localizer in a still more draconian way, we find equivalent localized categories. % [sic] et mêmes des équivalences -> et même des équivalences, adverb "même" printed with plural s

\msp{page 27}
That's rather strange - we'll see if it's really true!

\subsection*{III) Construction of $\Phi_0, \Psi_0 :$ \begin{tikzcd}[ampersand replacement=\&, baseline=(\tikzcdmatrixname-2-1.base), row sep=small, column sep=normal]
\& \und{Fibsc}_0S \\
\Cat/S \arrow[ur] \arrow[dr] \& \\
\& \und{Cofibsc}_0S.
\end{tikzcd}}

We must have
\[ \Phi_0(X) = \Phi r C(X) \qquad\qquad\qquad \Psi_0(X) = \Psi r C(X), \]
where
\[ r : \und{Parf}\,S \longrightarrow \Cat/S \]
is the canonical inclusion. As functors with values in \underline{Fibsc} resp. \underline{Cofibsc}, this already defines $\Phi_0$, $\Psi_0$ by
\[ \Phi_0 = \Phi r C \qquad\qquad\qquad \Psi_0 = \Psi r C, \]
\printedpage{17}
% head: A. Grothendieck : Dérivateurs, Ch. VI, Édition des Universités de Montpellier II et Paris VII
% foot: Version 9 novembre 2001
but we must see that $\Phi_0$, $\Psi_0$ take their values in the full subcategories with index 0. This is therefore equivalent to the following Lemma, applied to $Y = rC(X)$:

\textbf{Lemma.} \textit{Let $Y$ be perfect over $S$ (or merely a \textup{Hot}-fibered object over $S$). Then $\Phi(Y)$ is in $\und{Fibsc}_0S$, $\Psi(Y)$ in $\und{Cofibsc}_0S$.}

We must prove (for $\Psi(Y)$) that if $s \longrightarrow t$ is an arrow in $S$, then
\[ \Psi(Y)_s \longrightarrow \Psi(Y)_t \]
is a $W$-equivalence. Now this arrow identifies with
\[ Y/s \longrightarrow Y/t \]
and we know well that it is a Hot-equi\msp{page 28}valence, by definition of Hot-fibered objects, hence it is a $W$-equivalence.

We thus have the functors
\[
\begin{tikzcd}[column sep=large]
\und{Fibsc}_0S \arrow[r, shift left, "p"] & \Cat/S, \arrow[l, shift left, "\Phi_0"]
\end{tikzcd}
\]
we must see that they satisfy the general conditions of the lemma p.~11 (and it will be the same for $q$, $\Psi_0$). We already know that $p$ is compatible with the localizers. Let us prove it for $\Psi_0$. Let
\[ f : X \longrightarrow Y \]
in $\mathrm{Cat/S}$ be a $W$-equivalence, then $rC(f)$ is one too. It suffices to prove this: % [sic] Cat/S (upright S) -> Cat/$S$, the S printed upright unlike elsewhere

Let
\[ q : X' \longrightarrow Y' \]
be a morphism in $\und{Fibhot}\,S$ (for example a morphism of $\und{Parf}\,S$), which is a $W$-equivalence, hence \texttt{[which induces]} $W$-equivalences on the fibers. Then $\Psi(q) : \Psi(X') \longrightarrow \Psi(Y')$ is a $W$-equivalence fiber by fiber, i.e.
\[ X'/s \longrightarrow Y'/s \]
a $W$-equivalence for every $s \in S$ (was seen in corollary 2, page 20).

\msp{page 29}
Let us consider
\[ p\Phi_0 : \Cat/S \longrightarrow \Cat/S, \]
this thus transforms $X$ into $\Phi C(X)$, regarded as an object of $\Cat/S$ and nothing more. We have a functorial morphism
\[ X \xrightarrow{j_X} C(X) \qquad\qquad\qquad j : \mathrm{id}_{\Cat/S} \longrightarrow C, \]
\printedpage{18}
% head: A. Grothendieck : Dérivateurs, Ch. VI, Édition des Universités de Montpellier II et Paris VII
% foot: Version 9 novembre 2001
hence a functorial morphism
\[ \Phi(X) \xrightarrow{\Phi(j_X)} \Phi C(X) \qquad\qquad\qquad \Phi * j : \Phi \longrightarrow \Phi_0 \]
and also a functorial morphism
\[ X \xrightarrow{i_X} \Phi(X), \]
hence on composing
\[
\begin{tikzcd}
X \arrow[r, "i_X"] \arrow[rr, bend right, "\alpha_X"'] & \Phi(X) \arrow[r, "\Phi(j_X)"] & \Phi C(X),
\end{tikzcd}
\]
i.e.
\[ \alpha : \mathrm{id}_{\Cat/S} \longrightarrow p\Phi_0. \]
For $X$ fixed in $\Cat/S$, we know that $j_X$ is in $W$, hence so is $\Phi(j_X)$ (though not fiber by fiber), cf. cor.~1 page 20. We also know that $i_X$ is in $W$, hence so is the composite $\alpha_X$.

Let us consider
\[ \Phi_0 p : \und{Fibsc}_0\,S \longrightarrow \und{Fibsc}_0\,S\ \ . \]
We start with $X$ in $\und{Fibsc}_0\,S$, we forget the splitting and the very fact that it is Cat-fibrant
\msp{page 30}
(we have applied $p$), we apply $\Phi_0$ to it, i.e. we regard $\Phi CX$ as an object of $\und{Fibsc}\,S$ (and even of $\und{Fibsc}_0\,S$). We want again to define
\[ X \xrightarrow{\beta_X} \Phi_0 p(X) = \Phi C(X) \qquad\quad \text{\textbf{in} } \und{Fibsc}_0\text{\textbf{,}} \] % ?? uncertain: weight of "Fibsc_0" after bold "dans" -- regular or bold
functorial in $X$. Now we have as before
\[ X \xrightarrow{j_X} C(X), \]
hence on applying $\Phi$
\[ \Phi(X) \xrightarrow{\Phi(j_X)} \Phi C(X), \]
but regarded this time as an arrow in $\und{Fibsc}_0\,S$ (N.B. we do have $\Phi(X)$ in $\und{Fibsc}_0\,S$, for $X$, coming from $\und{Fibsc}_0\,S$, is in $\und{Fibhot}\,S$, and we apply the remark ``Cor.~3'' page 21.) On the other hand, we have an arrow in $\Cat/S$
\[ X \xrightarrow{i_X} \Phi(X), \]
is it an arrow in $\und{Fibsc}_0(S)$, i.e. is it cartesian? We see at once that no, that we have only, for an arrow $u : s \longrightarrow t$ in $S$, a commutation homomorphism (not an isomorphism) in

\end{document}
